\begin{summarybox}
	\textbf{\textcolor{CSummary}{一句话核心}}
	\textbf{塞瓦定理将“三线共点”的几何条件转化为面积比或线段比乘积为 $1$ 的代数方程，是证明共点与求比例的核心工具.}

	\medskip
	\textbf{\textcolor{CSummary}{使用条件}}
	\begin{itemize}[leftmargin=*, itemsep=0.5em, topsep=0.4em]
		\item \textbf{条件 1（点分边位）：} $D,E,F$ 分别位于 $\triangle ABC$ 的三边所在直线 $BC,CA,AB$ 上（可含延长线）.
		\item \textbf{条件 2（三线共点）：} 直线 $AD,BE,CF$ 交于同一点.
	\end{itemize}

	\medskip
	\textbf{\textcolor{CSummary}{关键提醒}}
	\begin{itemize}[leftmargin=*, itemsep=0.5em, topsep=0.4em]
		\item \textbf{易错点（有向符号）：} 各比值必须使用有向线段，点在延长线上对应比值为负，不能取绝对值.
		\item \textbf{检查项（顺序确认）：} 计算时严格按 $\frac{BD}{DC}\cdot\frac{CE}{EA}\cdot\frac{AF}{FB}$ 的顺序，避免颠倒.
	\end{itemize}

\end{summarybox}
